\subsection{Extended comparison ledger}

\begin{table}[H]
\centering
\scriptsize
\sisetup{round-mode=none, group-minimum-digits=20}
\setlength{\tabcolsep}{2pt}
\renewcommand{\arraystretch}{0.96}
\begin{tabularx}{\linewidth}{@{}>{\raggedright\arraybackslash}p{0.14\linewidth}>{\raggedright\arraybackslash}p{0.28\linewidth}>{\raggedright\arraybackslash}p{0.18\linewidth}>{\raggedleft\arraybackslash}p{0.15\linewidth}Y@{}}
\toprule
Row & TFPT source or target & \makecell[l]{Metrology\\readout} & \makecell[r]{Scheme\\representative} & Status note \\
\midrule
$m_\gamma$ & exact unbroken $U(1)_{\mathrm{em}}$ direction & \makecell[l]{$d_O=1$\\$m_\gamma/\lambda_\Sigma=0$} & $0$ & exact structural row \\
$m_g$ & exact color gauge direction before confinement & \makecell[l]{$d_O=1$\\$m_g/\lambda_\Sigma=0$} & $0$ & exact structural row \\
$W^\pm$ & threshold output at $M_Z$ & \makecell[l]{$d_O=1$\\$m_W/\lambda_\Sigma$} & \makecell[r]{\num{80.369}\\GeV} & scheme projection \\
$Z$ & threshold output at $M_Z$ & \makecell[l]{$d_O=1$\\$m_Z/\lambda_\Sigma$} & \makecell[r]{\num{91.188}\\GeV} & scheme projection \\
$H$ & threshold output & \makecell[l]{$d_O=1$\\$m_H/\lambda_\Sigma$} & \makecell[r]{\num{125.25}\\GeV} & scheme projection \\
\makecell[l]{$\hat s_Z^2$\\$(M_Z)$} & \makecell[l]{local gauge closure plus\\seam-even matching} & \makecell[l]{$d_O=0$\\scheme image} & \num{0.23116113} & scheme projection \\
\makecell[l]{$\alpha_s$\\$(M_Z)$} & \makecell[l]{local gauge closure plus\\singlet Wilson matching} & \makecell[l]{$d_O=0$\\scheme image} & \num{0.11796213} & scheme projection \\
$\Sigma m_\nu$ & intrinsic neutrino closure sum & \makecell[l]{$d_O=1$\\$\Sigma m_\nu/\lambda_\Sigma$} & \makecell[r]{\num{5.8764e-2}\\eV} & physical observable \\
$M_R$ & heavy-neutrino saddle on the closed neutrino / cosmology branch & \makecell[l]{$d_O=1$\\$M_R/\lambda_\Sigma$} & \makecell[r]{\num{1.31e15}\\GeV} & physical observable \\
$\Sigma m_\nu^{\mathrm{ss}}$ & appendix-level seesaw comparison from exact $\vgeo$ and $M_R$ & \makecell[l]{$d_O=1$\\$\Sigma m_\nu^{\mathrm{ss}}/\lambda_\Sigma$} & \makecell[r]{\num{5.779e-2}\\eV} & appendix-level auxiliary continuation \\
\bottomrule
\end{tabularx}
\caption{Supplementary gauge, Higgs, and neutrino rows for the extended comparison ledger. Exact
zero rows, physical observables, scheme projections, and appendix-level auxiliary continuations are
separated explicitly, with the dimensionless TFPT readout isolated from the chosen unitful
representative.}
\label{tab:extended-compiler-ledger}
\end{table}

\begin{table}[H]
\centering
\scriptsize
\begin{tabularx}{\linewidth}{@{}>{\raggedright\arraybackslash}p{0.11\linewidth}>{\centering\arraybackslash}p{0.16\linewidth}>{\raggedleft\arraybackslash}p{0.16\linewidth}Y@{}}
\toprule
Field & Word-length data $L_{f,j}$ & Derived transport matrix element & Status note \\
\midrule
$e$ & $8$ & \num{0.466} & Derived \\
$\mu$ & $5$ & \num{1.085} & Derived \\
$\tau$ & $3$ & \num{0.899} & Derived \\
$u$ & $7$ & \num{0.440} & Derived \\
$d$ & $7$ & \num{0.936} & Derived \\
$s$ & $5$ & \num{0.943} & Derived \\
$c$ & $3$ & \num{0.646} & Derived \\
$b$ & $2$ & \num{0.480} & Derived \\
$t$ & $0$ & \num{0.986} & Derived \\
\bottomrule
\end{tabularx}
\caption{Transport closure data as theorem outputs.}
\label{tab:yukawa-residuals}
\end{table}

\subsection{Novelty boundary and literature positioning}

The paper uses standard tools from relative spectral geometry, APS index theory, transport/resolvent language, and conventional RG comparison. The claim of novelty is not that these background ingredients are new in isolation, but that the minimal $3+2$ carrier packet, the compression identities, and the boundary-polarized closed architecture are organized into one tightly overdetermined program. Conversely, the manuscript does not claim novelty for standard benchmark conventions, nor does it claim that appendix-level information-theoretic or quasicrystal readings are part of the hard core.

\begin{thebibliography}{99}

\bibitem{aps1975}
M.~F.~Atiyah, V.~K.~Patodi, and I.~M.~Singer,
\newblock \emph{Spectral asymmetry and Riemannian geometry. I},
\newblock Math. Proc. Cambridge Philos. Soc. \textbf{77} (1975), 43--69.

\bibitem{chamseddineConnes1997}
A.~H.~Chamseddine and A.~Connes,
\newblock \emph{The spectral action principle},
\newblock Commun. Math. Phys. \textbf{186} (1997), 731--750.

\bibitem{codata2022}
P.~J.~Mohr, D.~B.~Newell, and B.~N.~Taylor,
\newblock \emph{CODATA recommended values of the fundamental physical constants: 2022 update},
\newblock NIST / CODATA reference set, \url{https://physics.nist.gov/cuu/Constants/}, accessed March 2026.

\bibitem{pdg2025}
S.~Navas et al.\ (Particle Data Group),
\newblock \emph{Review of Particle Physics},
\newblock Phys.\ Rev.\ D \textbf{110} (2024), 030001; 2025 online update, \url{https://pdg.lbl.gov/2025/}, accessed March 2026.

\bibitem{nufit2025}
NuFIT Collaboration,
\newblock \emph{NuFIT global analysis of neutrino oscillation data},
\newblock website snapshot based on data available through November 2025, \url{https://www.nu-fit.org/}, accessed March 2026.

\bibitem{planck2018}
Planck Collaboration,
\newblock \emph{Planck 2018 results. VI. Cosmological parameters},
\newblock Astron. Astrophys. \textbf{641} (2020), A6.

\bibitem{LiddleLeach2003}
A.~R.~Liddle and S.~M.~Leach,
\newblock \emph{How long before the end of inflation was observable inflation?},
\newblock Phys. Rev. D \textbf{68} (2003), 103503; arXiv:astro-ph/0305263.

\bibitem{Bombieri1966}
E.~Bombieri,
\newblock \emph{On exponential sums in finite fields},
\newblock Amer. J. Math. \textbf{88} (1966), 71--105.

\bibitem{DeligneWeil}
P.~Deligne,
\newblock \emph{La conjecture de Weil. I},
\newblock Inst. Hautes \'Etudes Sci. Publ. Math. \textbf{43} (1974), 273--307.

\bibitem{PolyaVinogradov}
G.~P\'olya and I.~M.~Vinogradov,
\newblock completion-of-sums lemma; classical references include
G.~P\'olya, \emph{\"Uber die Verteilung der quadratischen Reste und Nichtreste},
Nachr. K\"on. Ges. Wiss. G\"ottingen (1918), 21--29, and
I.~M.~Vinogradov, \emph{Sur la distribution des r\'esidus et des non-r\'esidus des puissances},
J. Phys.-Math. Soc. Perm \textbf{1} (1918), 94--98; modern textbook treatment:
H.~Iwaniec and E.~Kowalski, \emph{Analytic Number Theory},
AMS Colloquium Publications \textbf{53} (2004), \S 12.4.

\bibitem{minamiKomatsu2020}
Y.~Minami and E.~Komatsu,
\newblock \emph{New extraction of the cosmic birefringence from the Planck 2018 polarization data},
\newblock Phys. Rev. Lett. \textbf{125} (2020), 221301.

\bibitem{eskilt2023}
J.~R.~Eskilt et al.,
\newblock \emph{Cosmoglobe DR1 results. II. Constraints on isotropic cosmic birefringence from reprocessed WMAP and Planck LFI data},
\newblock Astron. Astrophys. \textbf{679} (2023), A144.

\bibitem{sullivan2025}
R.~M.~Sullivan, A.~Abghari, P.~Diego-Palazuelos, L.~T.~Hergt, and D.~Scott,
\newblock \emph{Planck PR4 (NPIPE) map-space cosmic birefringence},
\newblock arXiv:2502.07654.

\bibitem{na62obs2025}
NA62 Collaboration,
\newblock \emph{Observation of the $K^{+}\to\pi^{+}\nu\bar{\nu}$ decay and measurement of its branching ratio},
\newblock JHEP \textbf{02} (2025), 191; arXiv:2412.12015.

\bibitem{na62Moriond2026}
X.~Chang et al.\ (for the NA62 Collaboration),
\newblock \emph{New measurement of $K^{+}\to\pi^{+}\nu\bar{\nu}$ branching ratio at the NA62 experiment},
\newblock arXiv:2604.12649; contribution to the 2026 Electroweak session of the 60th Rencontres de Moriond.

\bibitem{koto2024}
KOTO Collaboration,
\newblock \emph{Search for the $K_{L}\to\pi^{0}\nu\bar{\nu}$ decay at the J-PARC KOTO experiment},
\newblock arXiv:2411.11237.

\end{thebibliography}
